%%% Fiktivní kapitola s ukázkami tabulek, obrázků a kódu

\chapter{Implementace návrhu}

Před zahájením vývoje aplikace je potřeba nainstalovat všechny potřebné nástroje a~nakonfigurovat vývojové prostředí. Nejdůležitější nástroje, které vyžadují globální instalaci, jsou následující:

\begin{itemize}
    \item Node.js
    \item Python
    \item Apache CouchDB
    \item Apache Solr
    \item Silk Workbench
    \item Apify CLI
\end{itemize}

V~řešení využijeme také řadu knihoven, které budeme instalovat pouze v~rámci projektu pomocí výchozího správce balíčků npm pro Node.js. Tyto knihovny vždy uvedeme na seznamu závislostí projektu, takže je lze snadno nainstalovat prostřednictvím příkazu \texttt{npm install}.

Co se týče výběru programovacích jazyků, přípravu dat vyřešíme pomocí vzájemně nezávislých skriptů psaných v~jazyce JavaScript. Samotnou aplikaci včetně serverové a~klientské vrstvy již napíšeme jazykem TypeScript, který je potřeba následně transpilovat do JavaScriptu. Tento dodatečný krok přidává komplexitu při spouštění kódu, proto jej vynecháme u~jednoduchých skriptů připravujících dokumenty pro databázi a~Solr. Zároveň ale přináší typovou kontrolu, kterou velmi oceníme v~komplexnější aplikaci a~to zejména při práci s~externími knihovnami, jejichž rozhraní není vždy perfektně zdokumentováno.

\section{Vývojové prostředí}

Aplikaci budeme vyvíjet v~editoru Visual Studio Code s~rozšířeními pro jazyky JavaScript a~TypeScript. Konzistentního formátování dosáhneme zapojením rozšíření Prettier, které má na starosti správné odsazení, maximální délku řádky a~další aspekty formátování.

Projekt je verzován ve vzdáleném repozitáři na platformě GitHub pod jménem MeaLinker\footnote{https://github.com/lhotanok/MeaLinker}. Název byl sestaven spojením slov \emph{meal} a \emph{linker}, snaží se totiž zachytit myšlenku propojení receptů z různých zdrojů. Dosavadní práce probíhala pro zjednodušení ve větvi \texttt{main}. Pokud by se na projektu začalo podílet více vývojářů, byl by zaveden tradiční systém vedlejších větví pro implementaci jednotlivých funkcí a požadavků na jejich sloučení s hlavní větví.

\section{Zpracování vstupních dat}

Jak jsme již zmínili v~části o~architektuře aplikace, data k~receptům získáme primárně pomocí vlastní extrakce dat. Pro dataset z~webové aplikace Food.com využijeme i~statická data z~platformy Kaggle, kterými rozšíříme extrahované informace. Data k~ingrediencím ze znalostních grafů DBpedia a~Wikidata získáme rovněž automatizovaným postupem skrze zpracování HTTP požadavků.

Jako zdroj receptů jsme si zvolili webové aplikace Food.com a~Allrecipes. Pro každou z nich vytvoříme dedikovaný extraktor s~využitím knihovny Apify. Nejprve se podíváme na řešení pro aplikaci Food.com, u~které nás na rozdíl od stránky Allrecipes čeká dodatečná fáze sloučení informací ze statického datasetu.

\subsection{Food.com}

Propojení dat se statickým datasetem Food.com Recipes and~Interactions nám do projektu zanese poměrně velkou komplexitu ve srovnání s~pouhou extrakcí dat z~webové aplikace, kterou využijeme u~stránky Allrecipes. Pro zpracování dat z~aplikace Food.com nemůžeme využít stejný přístup jako s~Allrecipes, protože neposkytuje plně strukturované ingredience (odděluje pouze množství od ostatního textu). Výrazná výhoda zpracování dat z~Food.com je v~jejich kvantitě, která s~více než $500\,000$ recepty $10\times$ převyšuje aplikaci Allrecipes a~tvoří polovinu obsahu největšího datasetu s~recepty Recipe1M+. Také poskytuje detailní informace přibližně k~$900$ ingrediencím na adrese \texttt{food.com/about}. Tyto data sice extrahovat nebudeme, počet ingrediencí s~detaily je pro nás ale užitečným vodítkem z~pohledu propojení s~otevřenými daty surovin.

Dataset Food.com Recipes and~Interactions obsahuje kolem $8000$ unikátních ingrediencí extrahovaných a~normalizovaných přímo z~jednotlivých receptů. Tyto normalizované ingredience budou po čištění velmi dobrým podkladem pro fasetový našeptávač ingrediencí na vyhledávací stránce. Využijeme je pro vyhledávání receptů z~libovolných zdrojů. Bez nich bychom našeptávač ingrediencí museli vytvořit ze jmen surovin z~jiných stránek, například Allrecipes. Případně bychom mohli využít externí seznam ingrediencí, ať už z~aplikace Food.com nebo z~otevřených znalostních grafů. Počet navrhovaných přísad by ale byl výrazně menší.

\subsubsection{Extrakce dat}

Všechny skripty pro přípravu dat z~Food.com jsou soustředěny v~adresáři \texttt{data/resources/foodcom}. V~této sekci budeme při uvádění cesty pro zjednodušení předpokládat, že se nacházíme uvnitř \texttt{data/resources/foodcom}. Samotná extrakce dat je umístěna v podadresáři \texttt{food-com-scraper}, jehož struktura odpovídá standardnímu Apify actorovi, kterého jsme představili v~předchozí kapitole. Pro vytvoření této struktury lze využít příkaz \texttt{apify\,create}, který uživatele provede možnostmi různých šablon. My zvolíme šablonu pro tzv. \texttt{CheerioCrawler}, což je v~knihovně Apify řešení využívající pouze HTTP požadavky a~parsování HTML z~odpovědi. Pokud bychom potřebovali automatizaci webového prohlížeče, použili bychom šablonu \texttt{PuppeteerCrawler}, případně \texttt{PlaywrightCrawler}. Šablony jsou pojmenovány podle klíčových technologií, tedy knihovny Cheerio pro extrakci dat z~HTML a~knihoven Puppeteer či Playwright poskytujících rozhraní pro automatizaci prohlížeče.

Pokud bychom adresář s~Apify actorem například naklonovali ze vzdáleného adresáře na platformě GitHub, stačilo by jej inicializovat příkazem \texttt{apify\,init}. Zároveň se jedná o~standardní Node.js projekt, pro instalaci potřebných knihoven tedy musíme spustit příkaz \texttt{npm\,install} nebo ekvivalentní \texttt{yarn\,install}. Po inicializaci se vytvoří důležitý adresář \texttt{apify\underline{{ }}storage} s~podadresáři \texttt{datasets} a~\texttt{key\underline{{ }}value\underline{{ }}stores}.

Do adresáře \texttt{datasets} se typicky ukládají JSON soubory s~jednotlivými položkami datasetu, v~našem případě recepty. Na cloudové platformě Apify se pak tyto soubory sloučí do společného datasetu, který lze stáhnout v~řadě různých formátů, z~nichž nejčastěji využívané jsou formáty JSON a~CSV. Naše řešení bude ale určeno pouze pro lokální vývoj, pro zjednodušení tedy budeme výsledky ukládat přímo sloučené do adresáře \texttt{key\underline{{ }}value\underline{{ }}stores/default} pod klíčem \texttt{RECIPES}. Zároveň zde pod klíčem \texttt{INPUT} najdeme další důležitý soubor se vstupem pro actora. Ten bude obsahovat objekt s~jedinou položkou \texttt{startUrls} v~následujícím formátu:

\begin{code}
{
  "startUrls": [
    {
      "url": "https://www.food.com/recipe/219662"
    },
    {
      "url": "https://www.food.com/recipe/103336"
    }
  ]
}
\end{code}
%$

Vstup bude vygenerován do \texttt{key\underline{{ }}value\underline{{ }}stores/default/INPUT.json} pomocí skriptu \texttt{recipe-urls-generator.js} z~adresáře \texttt{foodcom}, který URL adresy složí z~identifikátorů receptů. Před spuštěním tohoto skriptu je potřeba mít staženy soubory \texttt{recipes/RAW\underline{{ }}recipes.csv} a~\texttt{recipes/PP\underline{{ }}recipes.csv} z~datasetu Food.com Recipes and~Interactions\footnote{https://www.kaggle.com/datasets/shuyangli94/food-com-recipes-and-user-interactions} a~dokončen \texttt{recipes-preprocessing.js}. Přes proměnnou \texttt{RECIPES\underline{{ }}TO\underline{{ }}EXTRACT} uvnitř \texttt{constants.js} lze omezit počet receptů, které se mají extrahovat.

Kód Apify actora je typicky soustředěn v~adresáři \texttt{src}. Klíčové jsou soubory \texttt{main.js} s~inicializací (v našem případě třídy \texttt{CheerioCrawler}) a~\texttt{routes.js} s~funkcemi pro zpracování různých typů obrazovek. Obvykle je potřeba definovat chování pro stránku typu \texttt{LIST}, ze které jsou získány odkazy na detaily nalezených výsledků, a~pro stránku \texttt{DETAIL}, kde jsou extrahována a~uložena data jednotlivých výsledků. Naše řešení pro aplikaci Food.com zpracovává pouze detaily receptů, které obdrží na vstupu. Stačí tedy implementovat funkci \texttt{handleDetail}, která pomocí CSS selektoru \texttt{script[type="application/ld+json"]} zacílí JSON-LD data z~hlavičky HTML dokumentu, vybranou část z~nich převede do strukturované podoby a~nově zkonstruovaný recept přidá do souboru \texttt{RECIPES.json}.

Konstruktor třídy \texttt{CheerioCrawler} přijímá $1$ parametr s~možnostmi konfigurace. Prostřednictvím položky \texttt{proxyConfiguration} lze volitelně aktivovat rotování IP adres, které chrání naši vlastní IP adresu před dočasným či dokonce trvalým zablokováním a~zlepšuje poměr úspěšných požadavků. Vzhledem k~obecně nižší míře blokování ze strany aplikací s~recepty by využití proxy nemělo být nutné, je ale vhodné. Počet současně odesílaných požadavků omezíme na doporučenou hranici $50$ požadavků, abychom zamezili přetížení zpracovávané webové aplikace. Actora spustíme příkazem \texttt{apify\,run\,-p}. Příznak \texttt{-p} je zkratkou pro \emph{purge} a~zajišťuje korektní vyčištění předchozího stavu.

\subsubsection{Propojení dat}

V~první fázi potřebujeme extrahovat ingredience, které jsou poskytnuty v~souboru \texttt{ingr\underline{{ }}map.pkl}. Pro zpracování formátu \emph{pkl} potřebujeme knihovnu pickle určenou k~serializaci a~de-serializaci objektů jazyka Python \citep{pickle}. Je tedy potřeba aktivovat prostředí pro vývoj v~jazyce Python s~nainstalovaným balíčkem pickle (typicky virtuální prostředí označované jako \emph{venv}). Soubor převedeme skriptem \texttt{pkl-ingredients-extractor.py} do formátu CSV, kde nalezneme informace v~následujícím formátu:

\begin{code}
raw_ingr,raw_words,processed,len_proc,replaced,count,id
romaine lettuce leaf,3,romaine lettuce leaf,20,lettuce,4507,4308
iceberg lettuce leaf,3,iceberg lettuce leaf,20,lettuce,4507,4308
red romaine lettuce,3,red romaine lettuce,19,lettuce,4507,4308
\end{code}
%$

Klíčové pro nás budou unikátní hodnoty ze sloupců \texttt{replaced} a~\texttt{id}, přičemž každá hodnota \texttt{replaced} má přiřazeno unikátní \texttt{id}. Jména surovin normalizujeme na přísady pro vyhledávání pomocí skriptu \texttt{ingredients-preprocessing.js}. Dále z~nich vytvoříme jednoduchý RDF dataset ve formátu Turtle prostřednictvím skriptu \texttt{rdf-data/rdf-ingredients-converter.js} a~tento dataset nahrajeme do aplikace Silk Workbench pro nalezení linků s~grafy DBpedia a~Wikidata.

Normalizované názvy ingrediencí přidáme do receptů extrahovaných Apify actorem, což provedeme na základě informací ze souboru \texttt{RAW\underline{{ }}recipes.csv}, kde jsou propojeny recepty s~identifikátory ingrediencí. Stejným způsobem přidáme informace o~autorovi včetně jeho id. Tuto fázi má na starosti hlavní skript pro slučování dat, totiž \texttt{recipes-ingredients-merge-manager.js}.

Poslední fáze patří skriptu \texttt{ingredients-postprocessing.js}, který v~adresáři \texttt{rdf-data} očekává soubor \texttt{dbpedia-ingredients.json} s~extrahovanými ingrediencemi z~grafu DBpedia a~soubor \texttt{wikidata-ingredients.json} s~přísadami z~Wikidata. Ty generuje skript \texttt{rdf-data/external-dataset-linker.js}, který potřebuje soubory ve formátu N-triples \texttt{rdf-data/food-dbpedia-same-ingr.nt} a~\texttt{rdf-data/food-wikidata-same-ingr.nt} vytvořené pomocí Silk Workbench. Konfigurace linkování extrahovaná z~grafického rozhraní Silk Workbench je uložena v XML souborech ve stejném adresáři.

Výstupem našeho snažení jsou následující soubory:

\begin{code}
recipes/extended_recipes.json
ingredients/extended_ingredients.json
ingredients/search_ingredients.json
\end{code}
%$

První $2$~uvedené soubory jsou určeny pro uložení do databáze CouchDB a~poslední soubor najde uplatnění při tvorbě fasetových ingrediencí receptů uložených v~Solr. Skripty pro jednotlivé fáze zpracování jsou soustředěny ve skriptech \texttt{run} a \texttt{run-venv} vyžadující aktivované prostředí pro vývoj v Pythonu.

\subsection{Allrecipes}

\subsubsection{Extrakce dat}

Základní struktura actora pro extrakci dat z webové aplikace Allrecipes bude podobná jako u Food.com v předchozí sekci. Speciálně zpracování stránek s detaily receptů bude téměř identické, pouze využijeme dodatečné informace k ingrediencím obsažené v HTML elementech a uložíme je na pozici strukturovaných ingrediencí. Mezi actory ale bude podstatný rozdíl v získání odkazů na jednotlivé recepty. Actor pro Allrecipes nebude URL adresy přijímat na vstupu a namísto toho si je vyrobí sám. Má více způsobů, jak k extrakci přistoupit. My se podíváme na extrakci prostřednictvím interního API z vyhledávací stránky receptů a na typický přístup procházení jednotlivých kategorií.

\paragraph{Interní API}\mbox{}\\

Při vyhledávání receptů na základě filtrů můžeme přes vývojářské nástroje webového prohlížeče odchytit požadavek v následujícím formátu:

\begin{code}
/element-api/content-proxy/faceted-searches-load-more?page=1
\end{code}
%$
Formát odpovědi na tento požadavek je následující:
\begin{code}
{
    hasNext: boolean
    html: string
    totalResults: number
}
\end{code}
%$

Ze znalosti hodnoty {totalResults} a počtu výsledků na 1 stránce bychom měli být schopni předem určit, kolik stran budeme procházet. V době psaní této práce je inzerovaný počet výsledků uložený v položce \texttt{totalResults} roven $55\,680$. Poslední stranou při číslování od $1$~by tedy měla být strana $2\,320$. Při manuální kontrole lze ale zjistit, že poslední výsledky jsou na straně $417$, kde je zároveň uloženo \texttt{hasNext:\,false}.

Lze tedy předpokládat, že pokud na Allrecipes zadáme prázdnou množinu filtrů a prolistujeme všechny výsledky, neuvidíme inzerovaných $55\,680$ výsledků, ale pouze $(416 \cdot 24) + 16 = 10\,000$ receptů. Při podrobnějším zkoumání vypozorujeme, že aplikace Allrecipes používá maximální limit $10\,000$ výsledků i při zadání filtrů (například po přidání soli jakožto vyhledávací ingredience se celkový počet výsledků snížil na necelých $40\,000$ a poslední stranou výsledků je opět strana $417$. Pokud se tedy nespokojíme s $10\,000$ recepty, budeme muset přistoupit na řešení procházející jednotlivé kategorie receptů.

\paragraph{Kategorie receptů}\mbox{}\\

Odkazy na kategorie je možné získat ze stránky s relativní adresou \texttt{/recipes} přes CSS selektor třídy \texttt{recipeCarousel\underline{{ }{ }}link}. Kategorie pak můžeme procházet po stranách, stačí přidat query parametr \texttt{page}. Příklad relativní adresy pro $10$. stránku kategorie hlavních chodů: \texttt{/recipes/80/main-dish/?page=10}. Tento formát lze objevit při prozkoumání atributů tlačítka \texttt{Load\,more} na domovské stránce kategorie. Zda jsme již dorazili na poslední stranu poznáme podle obsahu elementu \texttt{h1}, ve kterém se na stránce bez receptů objeví hlášení \texttt{Page\,Not\,Found}.

Stránky kategorií budou v kontextu Apify actora odpovídat stránkám typu \texttt{LIST}. V rámci zpracování stránky kategorie musíme do fronty požadavků zařadit odkaz na další stranu (pokud již nejsme na poslední straně bez výsledků) a zároveň URL adresy detailů receptů. Fronta je oboustranná, můžeme tedy adresy s detaily receptů zařadit na začátek, aby se zpracovaly přednostně a měli jsme dříve k dispozici výsledky, neboť se průběžně ukládají do souboru \texttt{key\underline{{ }}value\underline{{ }}stores/default/RECIPES.json}.

Během extrakce dat ze stránek detailů receptů se musíme vypořádat s několika problémy. Aplikace Allrecipes není zcela konzistentní v mapování ingrediencí na atributy v HTML a střídá $2$ atributy pro jméno ingredience, které vidí uživatel, a název ingredience pro vyhledávání. Musíme tedy explicitně zkontrolovat, že pod položkou \texttt{text} ukládáme skutečně text, který je zobrazen bezprostředně po jednotce měření. Také narazíme na část receptů bez uvedených nutričních hodnot. Tyto recepty přeskočíme, neboť je pro nás informace o nutričních hodnotách důležitá.

\subsubsection{Propojení dat}

U datasetu z Allrecipes máme zjednodušenou práci, neboť strukturované ingredience máme již k dispozici z fáze extrakce receptů. Stačí nám tedy získat unikátní jména ingrediencí z extrahovaných receptů, z těchto ingrediencí vytvořit RDF dataset pro linkování s grafy DBpedia a Wikidata a extrahované informace propojit s dokumenty receptů. Postup je analogický jako s unikátními ingrediencemi z Food.com datasetu. 

Ingrediencím potřebujeme přiřadit identifikátor a ideálně se budeme snažit vyhnout duplicitě s ingrediencemi z Food.com, až je budeme ukládat do databáze spolu s rozšířením z DBpedia a Wikidata. Vygenerujeme si tedy uuid pro celou naši aplikaci, které budeme využívat jako tzv. \emph{seed} pro vytvoření dalších identifikátorů. Zde se nám bude hodit knihovna \texttt{uuid}\footnote{https://www.npmjs.com/package/uuid} a její funkce \texttt{v5}:
\begin{code}
uuid.v5(name, NAMESPACE_UUID)}
\end{code}
%$
Pokud takto vytvoříme identifikátory surovin stejného jména, namapují se na stejné uuid a vyhneme se duplicitě. Jména pro tento účel vždy převedeme do lowercase formátu.

\subsection{Linkování datasetů}

Pro nalezení linků mezi ingrediencemi z našich datasetů a entitami z DBpedia a Wikidata využijeme nástroj Silk Workbench, pomocí kterého vytvoříme RDF tvrzení s~IRI adresami ingrediencí spojenými vztahem \texttt{owl:sameAs}. V~rámci úlohy linkování navrhneme transformaci textu ingrediencí, která dokáže názvy propojit i~s~mírnými odlišnostmi ve formátu, čísle nebo pádu slov. Grafické znázornění jednotlivých fází transformace viz \ref{obr0a:silk-workbench}, příslušné XML soubory pro spuštění přes příkazovou řádku jsou u Food.com i Allrecipes v adresáři \texttt{rdf-data}.

\subsection{DBpedia}

Podklady pro extrakci dat ze znalostního grafu DBpedia jsou umístěny v adresáři \texttt{data/resources/dbpedia}. Přiložený skript není určen pro samostatné spouštění, pouze poskytuje rozhraní pro skripty konkrétních datasetů uvnitř podadresářů \texttt{rdf-data}. Před zahájením extrakce vždy nejprve identifikujeme entity, ke kterým potřebujeme zjistit dodatečné informace. Jakmile máme IRI adresy připraveny, teoreticky bychom mohli vytvořit jeden společný SPARQL dotaz pro všechny entity a~ten odeslat na DBpedia SPARQL endpoint. Dotaz by ale v~závislosti na počtu nalezených spojení mohl skončit příliš dlouhý a~nenechal by prostor pro škálování.

Zvolíme tedy alternativní řešení --- dynamicky vytvoříme sadu dotazů stejného formátu, každý s~přibližně $20$ IRI adresami. Tyto dotazy zpracujeme postupně a~výsledky uložíme do společného JSON datasetu. Výsledky si navíc od SPARQL endpointu můžeme vyžádat v~řadě různých formátů. Pro naše účely bude nejpraktičtější formát JSON-LD, jehož obsah využijeme v~hlavičkách HTML dokumentů. Se SPARQL endpointem lze komunikovat přes grafické rozhraní ve webovém prohlížeči nebo prostřednictvím HTTP GET požadavků. Pro snadnější automatizaci procesu extrakce využijeme druhou možnost, kde obsah dotazu předáme na místě query parametru s~názvem \texttt{query}. Šablona dotazu konstruujícího RDF graf je umístěná v souboru \texttt{ingredients-dbpedia.sparql}. IRI adresy ingrediencí jsou vloženy pomocí regulárních výrazů do množiny tzv. \texttt{VALUES}. Dotaz ve zjednodušené podobě vypadá následovně:

\begin{code}
PREFIX rdfs: <http://www.w3.org/2000/01/rdf-schema#>
PREFIX dbo: <http://dbpedia.org/ontology/>
PREFIX dbp: <http://dbpedia.org/property/>

CONSTRUCT {
    ?ingredient rdfs:comment ?comment ;
                rdfs:label ?label ;
                dbo:thumbnail ?thumbnail ;
}
WHERE {
    VALUES ?ingredient { 
        ##INGREDIENTS## 
        # Replace this section with ingredient resources such as:
        # <http://dbpedia.org/resource/Avocado>
        ##INGREDIENTS##
    }

    ?ingredient rdfs:label ?label .
    FILTER (LANG(?label) = "en")
    
    OPTIONAL {
        ?ingredient rdfs:comment ?comment .
        FILTER (LANG(?comment) = "en")
    }

    # Image with description(s)
    OPTIONAL {?ingredient dbo:thumbnail ?thumbnail .}
}
\end{code}
%$

\subsection{Wikidata}

Proces extrakce dat z grafu Wikidata bude probíhat analogicky k~postupu pro DBpedia z předchozího odstavce. I~zde využijeme HTTP GET požadavky na SPARQL endpoint, kde prostřednictvím query parametrů předáme obsah dotazu a~požadovaný formát výsledku. Projekt Wikidata neposkytuje reprezentaci JSON-LD, vystačíme si ale s~běžným JSON formátem, který lze vyžádat přes hodnotu query parametru \texttt{format} nastavenou na \texttt{json}. Z~této reprezentace pak sami vytvoříme odpovídající JSON-LD formát, který je vhodný pro strukturovaná data v~hlavičce HTML dokumentu.

Související soubory jsou uloženy v destinaci: \texttt{data/resources/wikidata}. Obsažený skript je opět volán pouze externě na základě extrahovaných ingrediencí z datasetů Food.com a Allrecipes. 

\section{Databáze Apache CouchDB}


\section{Vyhledávání pomocí Apache Solr}


\section{Middleware}


\section{Single-page aplikace}

V návrhu architektury klientské části aplikace jsme rozhodli, že naše řešení sestrojíme s využitím knihovny React a jejích funkcionálních komponent v kombinaci se speciálními funkcemi zvanými Hooks. Nejprve si tedy představíme tyto klíčové stavební prvky. Následně se zaměříme na jejich zapojení v rámci aplikace s recepty.

\subsection{Funkcionální komponenta}

Jedná se o~jednoduchou funkci přijímající tzv.~\emph{props} na místě parametru a~s~návratovým typem základního elementu knihovny React, tedy \texttt{JSX.Element}. Tyto elementy mohou být zapsány pomocí \emph{JSX}, což je syntaktické rozšíření jazyka JavaScript. Připomíná šablonovací jazyk, neboť kombinuje syntaxi jazyka HTML s~kódem psaným v~JavaScriptu. React zajišťuje renderování JSX elementů do HTML dokumentu prostřednictvím rozhraní DOM \citep{jsx-intro}. Příklad jednoduchého JSX elementu, který se do HTML vygeneruje jako tag \texttt{h1} s~textem \texttt{Headline}:

\begin{code}
const element = <h1>Headline</h1>;
\end{code}
%$

Ekvivalentem pro vývoj v~jazyce TypeScript je formát \emph{TSX}, který využijeme v~naší aplikaci. Rozlišujeme čistě prezentační komponenty, které pouze renderují data přijatá přes parametr \texttt{props}, a~stavové komponenty způsobující vedlejší efekty v~podobě změny stavu aplikace. Dle konvence má každá komponenta, byť jednoduchá, nárok na vlastní soubor, který snadno poznáme podle koncovky \texttt{jsx}, respektive \texttt{tsx} při vývoji v TypeScriptu.

\subsection{Hooks}

Koncept Hooks byl poprvé uveden ve verzi React $16.8$. Jedná se o speciální funkce, které umožňují spravovat stav, životní cyklus a další vlastnosti funkcionálních komponent bez využití tříd. Pomáhají dělit komponenty do menších funkcí na základě souvisejících částí kódu, což v komponentách na bázi tříd často není snadné zařídit. Například uvnitř vestavěné funkce \texttt{componentDidUpdate} se mohou potkat zcela nesouvisející části kódu, které je potřeba vykonat v případě aktualizace komponenty. Aplikační logiku soustředěnou uvnitř Hooks lze využít ve více komponentách a také je snadné kód testovat nezávisle na komponentě \citep{react-hooks}.

Pro pojmenování Hooks existuje striktní jmenná konvence --- všechny Hooks musejí začínat slovem \texttt{use}. Příklady nejdůležitějších Hooks poskytovaných knihovnou React jsou \texttt{useState} a \texttt{useEffect} vztahující se po řadě ke stavu a životnímu cyklu komponenty. V rámci jedné komponenty je lze podle potřeby využít libovolně mnohokrát. Také je možné vytvářet vlastní Hooks. Stejně jako komponenty se ukládají do souborů s koncovkou \texttt{tsx} a React je správně identifikuje jako Hooks právě na základě jména s počátečním slovem \texttt{use}. Kromě správného pojmenování musí Hooks dodržovat následující omezení \citep{react-hooks-w3c}:
\begin{itemize}
    \item Hooks mohou být použity pouze uvnitř funkcionálních komponent. V běžných funkcích ani komponentách implementovaných jako třída nebudou fungovat.
    \item Hooks nelze volat ve vnořených funkcích komponenty.
    \item Hooks nemohou ukládat různý výsledek na základě podmínky.
\end{itemize}
 
\subsection{Struktura aplikace}

Knihovnu React lze s aplikací integrovat různými způsoby. Nejjednodušším řešením je přidání \texttt{<script>} tagu do HTML dokumentu, volitelně s dalším tagem \texttt{<script>} pro podporu JSX syntaxe. Tento způsob je vhodný pro aplikace, které primárně nejsou navrženy jako single-page a potřebují prvky knihovny React využít pouze pro vybrané dynamické komponenty. Náš projekt má být naopak na principu single-page aplikace založen, využijeme tedy doporučené řešení v podobě prostředí Create React App\footnote{https://www.npmjs.com/package/create-react-app}. Adresářovou strukturu a konfiguraci projektu vytvoříme následujícím příkazem:
\begin{code}
npx create-react-app frontend
\end{code}
%$
Aplikaci spustíme na výchozí adrese \texttt{localhost:3000} spuštěním příkazu \texttt{npm start} v kořenovém adresáři \texttt{frontend}. Ve vytvořeném projektu je nastavena funkce automatické aktualizace při uložení souboru, díky které okamžitě vidíme provedené změny v prohlížeči při každém uložení.

\subsubsection{Výchozí HTML dokument}

V podadresáři \texttt{public} nalezneme kořenový HTML dokument \texttt{index.html}. Ten nebude potřeba z naší strany příliš modifikovat, pouze přidáme odkaz na ikony z knihovny Material UI, abychom je mohli využívat v rámci aplikace a také nahradíme obsah tagu \texttt{<title>} názvem našeho projektu, tedy MeaLinker. V adresáři \texttt{public} je rovněž uložena ikona aplikace. Využijeme logo vygenerované webovou aplikací Logo Maker\footnote{https://express.adobe.com/express-apps/logo-maker/} z dílny Adobe Express, které zadáme klíčové slovo \emph{Recipe}. Pokud bychom potřebovali banner včetně jména aplikace, přidali bychom do konfigurace generování loga ještě název \emph{MeaLinker}. Jako podklad použijeme ikonu\footnote{https://thenounproject.com/icon/search-food-2835433/} dohledatelnou pod frází \emph{search food}. Vygenerovaný výsledek je znázorněn obrázkem \ref{obr03:mealinker-logo}. Při návrhu barevného schématu aplikace se stejně jako u loga budeme držet červených odstínů, které doplníme o prvky zelené jakožto sekundární barvy palety.

\begin{figure}[h!]\centering
\includegraphics[width=40mm]{../img/mealinker-logo}
\caption{Logo aplikace vygenerované pomocí nástroje Logo Maker.}
\label{obr03:mealinker-logo}
\end{figure}

\subsubsection{Výchozí bod aplikace}

V adresáři \texttt{src} nalezneme soubor \texttt{index.tsx}, kde je vytvořen kořen stromové struktury elementů \texttt{React.ReactNode}, který zajišťuje renderování těchto elementů pomocí metody \texttt{render}. Volání metody \texttt{render} odpovídá architektuře znázorněné v předchozí kapitole, viz obrázek \ref{obr02:react-app}:
\begin{code}
const root = createRoot(
    document.getElementById('root') as HTMLElement
);

root.render(
  <React.StrictMode>
    <BrowserRouter>
      <App />
    </BrowserRouter>
  </React.StrictMode>,
);
\end{code}
%$

\subsubsection{Binární soubory}

V adresáři \texttt{src/assets} soustředíme potřebné binární soubory. V současné verzi aplikace zde budou pouze obrázky, konkrétně ikony využívané napříč aplikací. Většinu získáme z kolekce volně dostupných ikon platformy Flaticon\footnote{https://www.flaticon.com/}. Je vyžadován odkaz na zdroj a tvůrce ikony, s čímž si poradíme zahrnutím odkazů uvnitř tagu \texttt{<footer>}.

\subsubsection{Funkcionální komponenty}

Adresář \texttt{src} kromě obsahuje také podadresáře \texttt{recipes} a \texttt{ingredients}. Oba mají podobnou strukturu a obsahují adresáře \texttt{pages}, \texttt{components} a texttt{types}. Adresář texttt{types} obsahuje definice typů entit, které jsou přijímány ze serverové části aplikace. Za povšimnutí stojí koncovka souborů typů, která je rozdílná od většiny ostatních souborů na frontendu, tedy texttt{ts} namísto texttt{tsx}. Jedná se totiž o běžné soubory jazyka TypeScript, nikoli o komponenty určené specificky pro knihovnu React psané syntaxí TSX. Jinak jednotlivé typy odpovídají definicím na backendu aplikace, nemusíme se jim tedy dále věnovat.

\paragraph{Komponenty obrazovek}\mbox{}\\




\paragraph{Stavební komponenty}\mbox{}\\